%% 05_limitations.tex -- Owner group: G5. Rewritten for the camera-ready (2026-10-01).
%% No number that Results already gives; the CFG macros are numbers_cr_extra.tex (\NcrCfg*).
\section{Limitations}
\label{sec:limits}
All results come from one small regime. The grid has \Nconfigs\ configurations at \Nseeds\ seeds, with
\NminN--\NmaxN\,M non-embedding parameters, plus one \NmaxNfour\,M budget. The grid models are
trained for 30k steps on \NtrainHours\,h of English speech from one corpus (Emilia-EN), with one
codec (Mimi) and one confidence sampler. The 90k and 180k runs make more passes over the same
data. The ratio shrinks with training and is smaller at the largest budget and under prompt
CFG (\S\ref{sec:trend}), so we claim only the ordering, and only in this regime. A deeper grid is needed to locate the WER-optimal
depth (\S\ref{sec:scope}).

All metrics are automatic: ASR word error, speaker-encoder similarity and UTMOS. We ran no
listening test, so we cannot say whether a similarity gain of \NcrMatchLod\ or \NcrXEWLod\ is
audible. We did not evaluate MaskGCT, a masked-diffusion TTS system, on our items. F5-TTS is a calibration anchor: it differs from our
models in data, scale, codec and vocoder. Target length comes from a corpus-wide character
rate; per-item rate matching gave a small mean gain that a per-item sign test did not confirm
(Table~\ref{tab:negative}).
The selector (WavLM-L) and the ECAPA scorer both use an ECAPA-TDNN speaker head \citep{ecapa}
trained for verification on VoxCeleb data \citep{voxceleb1,voxceleb2}, so they may share
biases. Swapping their roles keeps search ahead of refinement at matched NFE
(Table~\ref{tab:search}), but at NFE 128 after $3\times$ training the swapped interval includes
zero (App.~\ref{app:search}). We did not try GE2E, x-vector or WavLM-base+ as selectors, because
they scored only the candidate WavLM-L selected. The selection cost omits resampling and the
speaker head of WavLM-L (App.~\ref{app:search}).

% PREREGISTRATION-v1.5 sec. 4(c) (verdict NEITHER, artifacts-v1.5/cfg_gate.json) keeps this
% check out of the findings; it is reported here only, with its deltas (reviewers Cx9t, ws2F).
The main grid uses no classifier-free guidance (CFG) \citep{cfg}. A separate registered check
retrained \NcrCfgRuns\ runs with condition dropout and compared prompt CFG at $T{=}16$ with
unguided $T{=}32$ at matched NFE. With guidance the ordering still holds, but the ratio
falls from \NcrCfgRatioUnguided\ (the same runs without guidance) to \NcrCfgRatio\ (CI
\NcrCfgRatioci), so the factor depends on the sampler as well as on scale. Guidance raised
SIM-o by \NcrCfgDsim\ (CI \NcrCfgDsimci) at a cost of \NcrCfgDwer\ absolute WER (CI
\NcrCfgDwerci), more WER than the registered limit allows; we did not run search on these
models.

We registered four hypotheses about the fitted surface, two of them primary
(\S\ref{sec:scope}). The step-exponent test passed, but its sign depends on the error
coordinate. The exchange-rate test failed its AICc criterion with $\kappa$ at its bound, and
the shape test could not be decided because the width coefficient sits at its cap. The extrapolation test failed. The share of the reachable range, our main
measurement, was added later (protocol v1.4), after $\Delta\tau$ was found to depend on the
coordinate. It is post hoc, not a registered test (App.~\ref{app:prereg}).
